\paragraph{The marked cover and the free lower rows.}
Write \(\mathfrak M\) for the conjugate of the marked half level
and \(\mathfrak J\) for the auxiliary type ideal. The level used
for theta spanning is 
\[
  \mathfrak m=\mathfrak M\mathfrak J,
\]
 with
\((\mathfrak M,\mathfrak J)=1\) and
\(\mathfrak J\overline{\mathfrak J}=(p_d)\).
In the group of invertible \(r\)-frames modulo \(\mathfrak m\)
take the kernel of reduction modulo \(\mathfrak M\). Chinese
remaindering identifies it with the general linear group on the
auxiliary factor. The finite test is supported where the upper
\(r\) rows reduce to the prescribed frame orbit. Its value there
is the auxiliary character in that frame.

The character is
trivial on the special linear subgroup. Indeed the finite ring
is a product of local rings whose residue fields have more than
three elements. Row reduction generates the special linear group
by elementary matrices, and conjugation by a diagonal unit whose
difference from one is a unit expresses each elementary matrix
as a commutator. In rank one the determinant is the identity.

\begin{proposition}[Theta evaluations at the marked level]
\label{prop:marked-theta-family}
For the prescribed half-level marking, the free lower rows at
\(\mathfrak m=\mathfrak M\mathfrak J\) realize the component evaluations at the shifted weights of Paragraph~\ref{par:shifted-theta}.
On component \(k\) write \(a_k(T),b_k(T)\) for the two complex evaluations of the rule \(T\) in the respective component coordinates. They are proportional by one factor on the component,
\[
 \exists C_k\in\mathbb C^\times\quad\forall T,\qquad a_k(T)=C_k b_k(T).
\]
This also applies when the desired family has different sections on its components.
\end{proposition}

\begin{proof}
\proofstep{The perfect pairing at the marked level.}
The principal polarization gives a perfect level pairing, with
the chosen primitive root of unity specifying its values. Write \(\delta_k\) for the dual half marking on component
\(k\), and choose generators \(b_i\) modulo \(\mathfrak M\)
on the reference module. At a split place a matrix \(A\)
on one half of a level pairing determines the matrix on the
opposite half by
\[
                  A^{\mathsf T}B=1,\qquad B=A^{-\mathsf T}.
\]
Thus one chooses a local basis on one half and the dual basis
on the other. Complete these frames on the auxiliary factor.
The resulting open level conditions can be satisfied by one
rational unitary move.

For this approximation use the Cayley transform
\(C(X)=(1-X)(1+X)^{-1}\) on the open set of skew-Hermitian
matrices with \(1+X\) invertible. One checks from \(X^*=-X\)
that 
\[
  C(X)^*C(X)=1.
\]
 This is the open chart of unitary
matrices with \(1+C(X)\) invertible. Its intersection with
each of the prescribed local open sets is dense by the chart
and its translates. Approximate the local parameters
simultaneously in the rational vector space of skew-Hermitian
matrices. The resulting rational unitary move satisfies the
same open level conditions.
The resulting moves take the period lattices into the indicated
genus with the same prescribed residues of \(p_d\delta_k(b_i)\)
modulo \(\mathfrak M\).

If \(gL_k=L_{k'}\) with \(g\) in the indicated level
group, the marking congruences read
\[
        p_d\,g\delta_k(b_i)
                \equiv p_d\,\delta_{k'}(b_i)\pmod{\mathfrak M}
                       \qquad(1\le i\le r).
\]
Cancel \(p_d\) and use the generating property of the \(b_i\).
The two dual markings agree on the entire module modulo
\(\mathfrak M\). The perfect pairing recovers the original
markings modulo the full rational level. Conjugating back by
the two moves gives an isomorphism of the marked polarized
period data. The uniformization identifies their moduli components.
This proves the orbit separation used above.

For a stabilizing move the same comparison puts its pullback in
the marked component group. Its determinant is an integral
norm-one unit and hence a root of unity. The principal congruence
of rational level at least three forces this torsion determinant
to be one.
Writing \(N\ge3\) for the rational principal level, the stabilizer element \(g\) satisfies
\[
  \det g\in\mu(E),\quad \det g\equiv1\pmod{N}\quad\Longrightarrow\quad\det g=1 .
\]
 For the signed vacuum, choose a symplectic basis
adapted to the split lattice and its half level.

In these
coordinates the off-diagonal blocks of the period change vanish
modulo eight and the determinant of its lower right block is
one modulo eight. These congruences are the adapted-basis
comparison for a principal-level isometry of determinant one.
For the lower right block \(D\), the comparison gives
\[
       \det D\equiv1\pmod8,\qquad \chi_4(\det D)=1 .
\]
In the squared half-characteristic transformation the
Gaussian factor is exactly the signed-vacuum coordinate change.
After taking the \(n/2\)-th power, the resulting determinant
factor cancels the one in the frame change. Here \(n\) is even. Evaluate the transformation
at a point where the theta function is nonzero. The remaining
vacuum scalar is one.
\label{cor:slice-marked-theta-cover-1}

\proofstep{Stable-range and full-level comparisons.}
There is one further level comparison in applying the stable
expansion. Its test requires zero lower rows modulo
\(\mathfrak m\), whereas the arithmetic rule tests lower rows
only modulo \(\mathfrak M\). Choose \(\alpha\) with

\[
  (\alpha)=\mathfrak J\mathfrak a
\]
 where \(\mathfrak a\) avoids
the levels and the bad primes. Multiply the lower rows by
\(\overline\alpha/\alpha\) and leave the upper rows fixed.
This is a unitary column change in the polynomial coordinates.
Let \(A\) be the transformed source lattice and \(S=\mathcal O_E^n\).
Take an integer \(M\) divisible by twice the bad and level
moduli. The construction gives \(t\) prime to \(M\) with
\(t\mathfrak J A\subseteq S\) and \(t\overline{\mathfrak J}S
\subseteq A\). Patch the lattice as
\[
              \Lambda'=\{v: Mv\in A,\ tv\in S\}.
\]
At the primes of \(M\) it equals \(S\), and at the other primes
it equals \(A\), by the two inclusions. Thus it is self-dual
and has an admissible level. On a target with shifted lattice
\(L^\sharp\), expand the moved vacuum into translates for
\(\Lambda'\). Its translation classes are indexed by

\[
  (A\otimes L^\sharp)/[(A\otimes L^\sharp)\cap
(\Lambda'\otimes L^\sharp)].
\]
 Choose representatives by
Chinese remaindering with upper corrections zero modulo
\(\mathfrak m\) and lower corrections zero modulo
\(\mathfrak M\). The upper rows still select the original
character. Conversely, keep the upper rows of any admissible
free-row argument and set the other rows to zero. This determines
its strict-test class. Subtract its moved representative from
the argument. The difference is in \(A\otimes L^\sharp\) and
determines the translation class uniquely. The even representatives
give trivial composition phase. This proves equality of the
two finite sums up to vacuum normalization. The positive factor
from projection is

\[
  [A:A\cap\Lambda']^r[\Lambda':A\cap\Lambda']^{-r/2}.
\]
The column change also multiplies the normalized vacuum by
a scalar of modulus one. Both factors are common to the targets,
so the harmonic polynomial, transformed by the same unitary
change, absorbs their ratio.
\label{cor:slice-marked-theta-cover-2}

\proofstep{A common scalar relating the two theta normalizations on each component.}
Finally, on a unit frame of component \(k\), the two normalized
vacua at the source of a short rule differ by one scalar. The
finite translate and derivative formulas therefore agree up to
this scalar for every rule simultaneously, after removing
\(u_{\Lambda,N_\Lambda}\) from each arithmetic reading. Its
modulus is fixed on that component. Write \(A_R\) for the normalized arithmetic reading of
a rule \(R\) after removal of \(u_{\Lambda,N_\Lambda}\), and
\(T_R\) for its classical reading. The vacuum comparison gives

\[
  A_R T_S=A_S T_R
\]
 for any two rules, in their common flag
character. If some \(T_S\) is nonzero, work first on its
nonvanishing locus. There \(A_S/T_S\) is holomorphic with fixed
positive modulus, so it equals a constant \(C_k\) on an open
set. The identity theorem on the connected frame domain gives
\(A_S=C_kT_S\) everywhere. For each \(R\) the cross-product identity now gives
\[
 (A_R-C_kT_R)T_S=A_R T_S-A_S T_R=0 .
\]
Cancel \(T_S\) on its nonvanishing locus; the identity theorem extends
the resulting equality over the connected frame domain. If all classical readings vanish, the vacuum
comparison makes the arithmetic readings zero too and any
scalar of the prescribed modulus serves. This is the common
\(C_k\) used in extending an arbitrary component family.
\label{cor:slice-marked-theta-cover-3}
\end{proof}
